\subsection{Результаты работы программы}

Рассмотрим пять функций с различными аналитическими особенностями.

\begin{enumerate}
    \item \textit{Функция с разрывом второго рода:}
    \[
    f_1(x) = \frac{1}{x}, \quad x \in \mathbb{R} \setminus \{0\}.
    \]
    Имеет бесконечный разрыв в точке $x = 0$. Интеграл в смысле Римана не существует при $0 \in [a, b]$. Алгоритм должен вернуть ошибку.

    \item \textit{Функция с устранимым разрывом:}
    \[
    f_2(x) = \frac{\sin x}{x}, \quad \lim_{x \to 0} f_2(x) = 1.
    \]
    В точке $x = 0$ функция не определена, но предел существует. Алгоритм должен корректно вычислить интеграл, используя алгоритмическое среднее или односторонние пределы.

    \item \textit{Квадратный многочлен:}
    \[
    f_3(x) = x^2 + 2, \quad x \in \mathbb{R}.
    \]
    Не имеет особенностей.

    \item \textit{Линейная функция:}
    \[
    f_4(x) = 2x + 2, \quad x \in \mathbb{R}.
    \]
    Метод трапеций должен дать точное значение интеграла уже при первом разбиении.

    \item \textit{Натуральный логарифм:}
    \[
    f_5(x) = \ln x, \quad x > 0.
    \]
    Не определена при $x \le 0$. При $a \le 0$ алгоритм должен вернуть ошибку; при $a > 0$ интеграл существует и вычисляется численно.
\end{enumerate}

Ниже приведены результаты интегрирования указанных функций методом трапеций с точностью $\varepsilon = 10^{-5}$ и лимитом разбиений $M = 10^6$.

\begin{table}[H]
\centering
\begin{tabular}{|c|c|c|c|c|c|c|}
\hline
$i$ & $[a, b]$ & $I_{\text{exact}}$ & $I_{\text{num}}$ & $|I_{\text{num}} - I_{\text{exact}}|$ & Разбиения & Статус \\ \hline
1 & $[0.1, 2.0]$ & $2.99573$ & $2.99573$ & $1.1789 \cdot 10^{-7}$ & 16384 & \texttt{OK} \\ \hline
1 & $[-1.0, 1.0]$ & --- & --- & --- & 0 & \texttt{ERR} \\ \hline
2 & $[-1.0, 1.0]$ & $1.89217$ & $1.89217$ & $1.91478 \cdot 10^{-7}$ & 1024 & \texttt{WARN} \\ \hline
2 & $[0.0, 1.0]$ & $0.946083$ & $0.946083$ & $9.5739 \cdot 10^{-8}$ & 512 & \texttt{WARN} \\ \hline
3 & $[0.0, 3.0]$ & $15$ & $15$ & $2,68221\cdot 10^{-7}$ & 1 & \texttt{OK} \\ \hline
4 & $[0.0, 5.0]$ & $35$ & $35$ & $0$ & 1 & \texttt{OK} \\ \hline
5 & $[1.0, 2.71828]$ & $1$ & $1$ & $1.48323 \cdot 10^{-7}$ & 1024 & \texttt{OK} \\ \hline
5 & $[-1.0, 1.0]$ & --- & --- & --- & 0 & \texttt{ERR} \\ \hline
\end{tabular}
\caption{Сводная таблица результатов тестирования метода трапеций.}
\end{table}